% =====================================================================
%  effective-writing-handout.tex
%  Handout 2 -- Writing for Readers: Paragraph Analysis Criteria
%
%  Source: L. McEnerney, "Leadership Lab: Writing Beyond the Academy",
%          University of Chicago, 23 Jan 2015 (recorded lecture).
%  Pairs with Handout 1 (mcenerney-lens): that handout checks a paper's
%  problem and positioning; this one checks a single paragraph.
%
%  REPO NOTES
%  - Criterion IDs (EW-1 ... EW-10), principle IDs (P1 ... P7) and
%    ground-rule IDs (G1 ... G5) are stable. Skills and reports cite
%    them, so never renumber: append new IDs at the end instead.
%  - Time stamps (h:mm:ss) point to the recorded lecture so every
%    criterion can be traced back to where it is taught.
%  - Compiles with pdfLaTeX on Overleaf (TeX Live 2020 or later).
%    For A4 paper, change "letterpaper" in the geometry line.
%  - 25 Oct 2026, settled with Handout 1 for the skills repo: EW-3 judges
%    stakes as Handout 1's stakes field does (a cost to these readers), so Version
%    B's EW-3 is PARTIAL; G6 (someone else's paragraph) appended;
%    P1 S3 labels for several paragraphs; Mode line and the covered-by-
%    lens line in the report; the title block says what each handout
%    leaves to the other. In Version B, S3's "it drains" became "the
%    closure drains", because "it" first read as backscatter. EW-3's
%    Check names the readers' predictions, results or decisions.
%  - 6 Nov 2026: cross-references follow Handout 1's renumbered fields
%    (problem L4, stance L5, stakes L6, contribution L8, opening L12).
% =====================================================================

\documentclass[11pt]{article}

% ---------- Fonts and encoding ---------------------------------------
\usepackage[T1]{fontenc}
\usepackage[utf8]{inputenc}
\usepackage{charter}                 % body: Bitstream Charter
\usepackage[scaled=0.90]{helvet}     % sans: Helvetica clone for headings
\usepackage{microtype}

% ---------- Page layout ----------------------------------------------
\usepackage[letterpaper,margin=0.9in,headheight=14pt]{geometry}
\usepackage{xcolor}
\usepackage{enumitem}
\usepackage{array,tabularx,booktabs}
\usepackage{titlesec}
\usepackage{fancyhdr}
\usepackage[most]{tcolorbox}
\usepackage[hidelinks]{hyperref}
\hypersetup{pdftitle={Writing for Readers: Paragraph Analysis Criteria (Handout 2)}, pdfauthor={Mohammed Meligy}}

\setlength{\parindent}{0pt}
\setlength{\parskip}{5pt}

% ---------- Colours --------------------------------------------------
\definecolor{ink}{HTML}{1F2328}
\definecolor{accent}{HTML}{1F4E79}
\definecolor{muted}{HTML}{5B6B7D}
\definecolor{rulec}{HTML}{C9D3DF}
\definecolor{fill}{HTML}{F2F5F9}
\definecolor{strongc}{HTML}{2E6B3F}
\definecolor{partialc}{HTML}{9A6700}
\definecolor{missingc}{HTML}{A33A2B}
\color{ink}

% ---------- Headings -------------------------------------------------
\titleformat{\section}{\sffamily\large\bfseries\color{accent}}{\thesection}{0.6em}{}
\titleformat{\subsection}{\sffamily\normalsize\bfseries\color{ink}}{}{0em}{}
\titlespacing*{\section}{0pt}{16pt}{6pt}
\titlespacing*{\subsection}{0pt}{10pt}{3pt}

% ---------- Header and footer ----------------------------------------
\pagestyle{fancy}
\fancyhf{}
\renewcommand{\headrulewidth}{0.4pt}
\renewcommand{\headrule}{\color{rulec}\hrule height\headrulewidth}
\fancyhead[L]{\sffamily\footnotesize\color{muted}Writing for Readers \textperiodcentered{} Paragraph Analysis Criteria}
\fancyhead[R]{\sffamily\footnotesize\color{muted}effective-writing}
\fancyfoot[C]{\sffamily\footnotesize\color{muted}\thepage}

% ---------- Small helpers --------------------------------------------
\newcommand{\ts}[1]{{\sffamily\footnotesize\color{muted}[#1]}}        % lecture time stamp
\newcommand{\sn}[1]{\textsuperscript{\sffamily\bfseries\color{accent}S#1}\,} % sentence label in example text
\newcommand{\sref}[1]{{\sffamily\footnotesize\bfseries\color{accent}S#1}}     % sentence reference in running text
\newcommand{\srange}[2]{{\sffamily\footnotesize\bfseries\color{accent}S#1--S#2}}
\newcommand{\id}[1]{{\sffamily\bfseries\color{accent}#1}}              % ID in running text
\newcommand{\badge}[2]{{\setlength{\fboxsep}{1.6pt}\colorbox{#1}{\sffamily\scriptsize\bfseries\color{white}\strut#2}}}
\newcommand{\stS}{\badge{strongc}{MET}}
\newcommand{\stP}{\badge{partialc}{PARTLY MET}}
\newcommand{\stM}{\badge{missingc}{NOT MET}}
\newcommand{\stN}{\badge{muted}{DOES NOT APPLY}}

% Lists
\setlist{itemsep=2pt,parsep=0pt,topsep=3pt}
\newlist{principles}{description}{1}
\setlist[principles]{style=multiline,leftmargin=2.6em,font=\sffamily\bfseries\color{accent}}
\newlist{fields}{description}{1}
\setlist[fields]{style=multiline,leftmargin=5.4em,font=\sffamily\small\bfseries\color{muted},itemsep=4pt,topsep=2pt}

% Criterion card:  \begin{criterion}{ID}{Title}{time stamps} ... \end{criterion}
\newtcolorbox{criterion}[3]{%
  enhanced,
  colback=white, colframe=rulec, boxrule=0.6pt, arc=2pt,
  colbacktitle=fill, coltitle=ink, toptitle=3pt, bottomtitle=3pt,
  left=8pt, right=8pt, top=5pt, bottom=5pt,
  fonttitle=\sffamily\bfseries,
  title={\textcolor{accent}{#1}\quad #2\hfill{\mdseries\footnotesize\color{muted}#3}},
  before skip=9pt, after skip=9pt}

% Plain shaded panel for notes and examples
\newtcolorbox{panel}[1][]{%
  enhanced, breakable, colback=fill, colframe=fill, arc=2pt,
  left=10pt, right=10pt, top=7pt, bottom=7pt,
  before skip=8pt, after skip=8pt, #1}

% =====================================================================
\begin{document}

% ---------- Title block ----------------------------------------------
{\sffamily\footnotesize\bfseries\color{muted}HANDOUT 2 \textperiodcentered{} EFFECTIVE-WRITING}\par\vspace{2pt}
{\sffamily\LARGE\bfseries Writing for Readers: Paragraph Analysis Criteria}\par\vspace{3pt}
{\sffamily\small\bfseries Mohammed Meligy}{\sffamily\small\color{muted}\ \textperiodcentered{} Mechanical and Aerospace Engineering, Arizona State University \textperiodcentered{} \textcopyright{} 2026, CC BY 4.0}\par\vspace{3pt}
{\small\color{muted}Compiled from \emph{Writing Beyond the Academy}, a Leadership Lab lecture by Larry McEnerney
(University of Chicago Writing Program), 23 January 2015. Unofficial study handout: the lecture's ideas are
paraphrased, and time stamps such as \ts{1:08:46} point to the recording. Pairs with Handout~1
(\emph{mcenerney-lens}), which checks a whole paper; used together on one draft, Handout~1 owns the opening's
problem, stance, stakes and contribution (L4, L5, L6, L8, L12) and this handout checks only its sentences
(\id{EW-5}--\id{EW-10}).}\par
\vspace{4pt}{\color{rulec}\rule{\linewidth}{0.6pt}}

\begin{panel}
\textbf{How to use this handout.}
\begin{enumerate}[leftmargin=1.6em]
  \item Fill in the \emph{reader card} (Section~2). No criterion can be judged without it.
  \item Number the sentences of your paragraph \sref{1}, \sref{2}, \sref{3}, \ldots{} For several paragraphs,
        label them too: P1~S1, P2~S3.
  \item Run criteria \id{EW-1} to \id{EW-10} in order. Give each a status and point to the sentences behind it.
        If you already ran Handout~1 on this draft, mark \id{EW-1}--\id{EW-4} for its opening paragraph
        ``covered by the lens'' and run \id{EW-5}--\id{EW-10} there.
  \item Fix the one criterion that most limits value first, usually \id{EW-1}--\id{EW-3}, before polishing sentences.
\end{enumerate}
\textbf{Someone else's paragraph:} say whether each criterion is met and what the paragraph does, with sentence
numbers; give no fix move (G6).\par\smallskip
\textbf{Status key:}\enspace \stS{}\quad \stP{}\quad \stM{}\quad \stN{}\quad (Handout~1: STRONG, PARTIAL, MISSING, N/A)
\end{panel}

% =====================================================================
\section{The idea behind the criteria}

In school, readers were paid to read you: teachers read your writing to assess what you know. After school,
nobody is paid to care about you. Readers read only while the text is valuable to them, and they stop when it
isn't. The job of a text you send is to change what its readers think, do or decide, and the readers are the ones
who judge whether it did. Every criterion below follows from that. \ts{11:21--18:50}

\begin{principles}
  \item[P1] \textbf{Value is the quality that matters.} Clear, organized and persuasive writing serves value; it
        does not replace it. \ts{5:50}
  \item[P2] \textbf{No rule about a text makes sense until you say who reads it and what it is for.}
        ``Use short sentences,'' ``avoid the passive,'' ``don't start with \emph{And},'' ``no jargon'' are
        techniques, not rules. Spelling is the one exception. \ts{7:26; 34:07}
  \item[P3] \textbf{Language is a technique for controlling the reader's reading process.} \ts{37:53}
  \item[P4] \textbf{Readers read in time.} When the text gets in the way they slow down, then misunderstand, then
        get annoyed, then stop. \ts{28:35}
  \item[P5] \textbf{Readers don't trust you.} They expect you to waste their time, so each sentence has to earn
        the next one. \ts{1:04:45}
  \item[P6] \textbf{Change their heads; don't reveal yours.} School trained you to demonstrate knowledge. That
        habit is the main obstacle. \ts{1:01:45; 1:15:45}
  \item[P7] \textbf{Think by writing, then revise for readers.} Drafting to discover your ideas is good. The
        criteria apply to the text you send, not to the draft you think with. \ts{9:54}
\end{principles}

% =====================================================================
\section{Before you analyze: the reader card}

\renewcommand{\arraystretch}{1.35}
\begin{tabularx}{\linewidth}{@{}>{\sffamily\small\bfseries}p{3.4cm}X@{}}
\toprule
Readers            & Who, specifically, will read this? A journal's community, a reviewer, a committee, a client. \\
Current belief     & What do they already think, assume or do that this paragraph touches? \\
What they value    & What would make reading this worth their time? \\
Function           & After this paragraph, what should they think, do or decide differently? \\
Reading conditions & At a desk or skimming? Will they get to the end of a long sentence? \\
\bottomrule
\end{tabularx}
\renewcommand{\arraystretch}{1}

\subsection{Ground rules for whoever analyzes (writer, colleague or AI)}
\begin{principles}
  \item[G1] \textbf{No finding without the reader card.} If the readers or the function are unknown, ask, or
        state the assumption you used.
  \item[G2] \textbf{Never report a text rule as a finding on its own.} Translate it into a reader effect: not
        ``uses the passive voice'' but ``the subject of \sref{3} is the apparatus, while these readers care about
        what it measured.''
  \item[G3] \textbf{Every finding cites sentence numbers.}
  \item[G4] \textbf{Judge concision by the reader's processing time, not by word count.}
  \item[G5] \textbf{Diagnose first; suggest one fix move per criterion.} Keep the writer's content, terminology,
        numbers and technical claims intact.
  \item[G6] \textbf{Someone else's paragraph gets no fixes.} When the text is not the user's, say whether each
        criterion is met and what the paragraph does, with sentence numbers, but give no fix move and no priority
        fix: it isn't the user's to revise.
\end{principles}

% =====================================================================
\section{The criteria}

\subsection{\texorpdfstring{A.\enspace}{A. }Value: does the paragraph give readers a reason to read?}

\enlargethispage{2\baselineskip}% room for the author line on page 1, so EW-1 stays on page 2
\begin{criterion}{EW-1}{The reader's role in the opening}{39:48--47:37; 1:03:04}
\begin{fields}
  \item[Ask] Does the first sentence give the reader a role, or only announce a topic?
  \item[Check] Place \sref{1} on the opening ladder:\\[2pt]
    \begin{tabular}{@{}>{\sffamily\bfseries\color{accent}}l l@{}}
      L0 & says what the text is \emph{about} (the reader has no role) \\
      L1 & says what the text \emph{argues} \\
      L2 & says what \emph{question} the text answers \\
      L3 & names the question the \emph{readers} have, or ought to have \\
    \end{tabular}
  \item[Red flags] ``This paper/section/paragraph discusses, outlines, presents, describes\ldots''; opening on
    background or definitions; the writer as the only actor in the sentence.
  \item[Fix] Rewrite \sref{1} so it starts from the readers' situation: what they think, what they need, or what
    they asked.
  \item[Applies] Any paragraph; essential for the first paragraph of a paper, section or memo.
\end{fields}
\end{criterion}

\begin{criterion}{EW-2}{The readers' belief, named and challenged}{49:11--55:12}
\begin{fields}
  \item[Ask] Does the paragraph say what the readers currently think, and why that is incomplete or wrong?
    Outside the academy the move is often different: say what the readers need, and that this text supplies it.
  \item[Check] Find the sentence that states the community's assumption or practice. Then find the turn:
    \emph{however, but, unrealistic, it is claimed, incomplete, inconsistent}.
  \item[Red flags] The paragraph describes the world, not the readers' understanding of it. A turn word with
    nothing before it to turn against.
  \item[Fix] Add one sentence that states the readers' view fairly, then make the turn.
  \item[Note] Most academic writing makes this move, usually politely. McEnerney's estimate is 85--95\% of it:
    in effect, ``I know what you think, and you're wrong.''
\end{fields}
\end{criterion}

\begin{criterion}{EW-3}{Stakes and urgency}{55:12--1:00:02}
\begin{fields}
  \item[Ask] Can these readers tell what it costs \emph{them} if they ignore this, or what they gain if they act
    on it?
  \item[Check] Name the cost in one phrase, as a cost to the readers on the reader card. A consequence stated
    only as a fact about the world (a physical effect, a number) is only partly met until it is tied to something these
    readers care about, such as their predictions, results or decisions. Then circle the words that carry tension (\emph{fails, cannot, error, risk, crisis}).
  \item[Red flags] Neutral placeholders: \emph{situation, aspects, implications, issues, has been studied}. A
    paragraph in which nothing is at risk.
  \item[Fix] Replace the neutral noun with the concrete cost: what breaks, for whom, and how badly.
  \item[Note] Readers are drawn to instability and trouble. ``Yesterday, nothing changed'' is not a headline.
    Handout~1 judges stakes the same way (L6).
\end{fields}
\end{criterion}

\begin{criterion}{EW-4}{Words that signal value}{48:05--49:04}
\begin{fields}
  \item[Ask] Do the words that describe your contribution tell readers they can use it?
  \item[Check] Circle every noun that names what you offer. Sort them into \emph{use} words (tool, method,
    test, criterion, diagnostic, a model that predicts\ldots) and \emph{display} words (ideas, concepts,
    thoughts, discussion, overview, insights).
  \item[Red flags] The contribution is described only with display words, or only as an activity (``X is
    analyzed'').
  \item[Fix] Name what the reader can do with it.
\end{fields}
\end{criterion}

\subsection{\texorpdfstring{B.\enspace}{B. }Sentences: is the value where readers look for it?}

\begin{criterion}{EW-5}{Focus: subjects carry what readers care about}{1:05:45--1:15:45}
\begin{fields}
  \item[Ask] Are the grammatical subjects the things these readers want to think about?
  \item[Check] Underline the subject of each sentence; in English it is the focus about four times in five.
    List the subjects in order. Do the readers care about each one? Do the subjects stay on one line of interest?
  \item[Red flags] Subjects that are the writer's apparatus (\emph{this study, the authors, the code,
    Section~3}) when readers care about the phenomenon. Subjects that jump between unrelated things.
  \item[Fix] Move what the readers care about into subject position. ``The dog chased the cat'' and ``The cat
    was chased by the dog'' are equally good until you know whether the readers are thinking about the cat.
  \item[Note] The same opening can focus on \emph{studies of} a topic, on \emph{the topic itself}, or on
    \emph{its readers}; McEnerney shows three such versions of one paragraph. Choose by readers, not taste.
\end{fields}
\end{criterion}

\begin{criterion}{EW-6}{Stress: endings carry what matters}{30:35--33:06; 1:08:46}
\begin{fields}
  \item[Ask] Do the ends of sentences hold the information you most want readers to notice?
  \item[Check] Circle the last few words of each sentence and read only the circled words, in order. Is that the
    paragraph's message?
  \item[Red flags] Endings that hold hedges, citations, locations (\emph{in this context, as shown in Fig.~3})
    or old information. A run of short sentences that stresses everything equally, so nothing stands out.
  \item[Fix] Move new, important information to the end. Use longer sentences with lesser stresses that build
    toward the main one.
\end{fields}
\end{criterion}

\begin{criterion}{EW-7}{Value along the way}{29:29--30:29}
\begin{fields}
  \item[Ask] In a long sentence, does the reader get something before the end?
  \item[Check] For each long sentence, mark the first point where the reader gains something. If the payoff
    comes only at the end, ask whether these readers will get there (reader card: reading conditions).
  \item[Red flags] Several lines of setup before the point, in a venue read quickly or on screen.
  \item[Fix] Bring an early payoff forward, or split the sentence. Shortness is not the goal: a long sentence
    that pays off as it goes is fine anywhere.
  \item[Note] Readers of a specialist journal will often reach the end of a long sentence; general readers
    often won't.
\end{fields}
\end{criterion}

\begin{criterion}{EW-8}{Concision as processing time}{1:06:46--1:08:24}
\begin{fields}
  \item[Ask] Will the reader have to process any sentence twice?
  \item[Check] Mark sentences where the reader must re-sort who did what, hold the subject across a long
    interruption before the verb, or flip the focus to the thing they care about.
  \item[Red flags] Active or passive chosen by rule rather than by focus. Nominalizations that hide the actor.
    A short sentence that still forces a reread.
  \item[Fix] Rebuild the sentence around the reader's focus, even if it gets a few words longer.
\end{fields}
\end{criterion}

\subsection{\texorpdfstring{C.\enspace}{C. }Reading process: does the paragraph keep readers reading?}

\begin{criterion}{EW-9}{Reading friction}{27:43--29:29; 6:43}
\begin{fields}
  \item[Ask] Where would a target reader first slow down?
  \item[Check] Mark the first friction point and name its cause: a term these readers don't use, jargon that
    adds nothing for them, an ambiguity, or a missing link between sentences.
  \item[Red flags] ``It's the reader's job to read it.'' Jargon judged by a rule rather than by the readers.
  \item[Fix] Fix the first friction point first. Keep the jargon these readers use and value; explain or cut the
    rest.
  \item[Note] A friction-free paragraph can still have no value. Ease is not the goal; it only keeps readers
    reading long enough to find the value.
\end{fields}
\end{criterion}

\begin{criterion}{EW-10}{Changing versus demonstrating}{11:21--12:25; 1:01:45}
\begin{fields}
  \item[Ask] Does each sentence change what the readers think, or only show what the writer knows?
  \item[Check] Tag each sentence \textbf{C} (moves the reader's thinking or decision) or \textbf{D}
    (demonstrates knowledge, method or effort).
  \item[Red flags] A paragraph of D sentences. Method recital or background with no consequence for the reader.
    ``It is well known that\ldots'' with nothing done to it.
  \item[Fix] Connect each D sentence to a change (why should these readers care that you know this?), or cut it.
\end{fields}
\end{criterion}

% =====================================================================
\clearpage
\section{Worked example}

\textbf{Reader card.} \emph{Readers:} LES modelers who use eddy-viscosity closures. \emph{Current belief:} the
subgrid model's job is to drain energy from the resolved scales. \emph{Value:} knowing where their closure will
fail. \emph{Function:} test the closure in the regions where it is weakest.

\begin{panel}[breakable=false]
\textbf{Version A}\quad
\sn{1}This section discusses subgrid-scale modeling in large-eddy simulation (LES) of turbulent flows.
\sn{2}Subgrid-scale models represent the effect of the unresolved scales on the resolved field.
\sn{3}The Smagorinsky model, which relates the subgrid stress to the resolved strain rate through an eddy
viscosity, is widely used.
\sn{4}In this work, the energy transfer between resolved and subgrid scales is analyzed using filtered DNS data.
\end{panel}

\begin{panel}[breakable=false]
\textbf{Version B}\quad
\sn{1}Most LES practitioners close the subgrid stress with an eddy viscosity, and in doing so they assume that
energy flows only from the resolved scales to the subgrid scales.
\sn{2}A priori tests on filtered DNS data show that this assumption is incomplete: locally, energy also flows
back to the resolved scales.
\sn{3}Because a positive eddy viscosity is purely dissipative, the model cannot represent this backscatter, and
where backscatter is strong the closure drains energy from scales that are in fact gaining it.
\sn{4}This paper gives modelers a tool for locating those regions, so that they can test their closure where it
is most likely to fail.
\end{panel}

{\footnotesize
\renewcommand{\arraystretch}{1.1}
\begin{tabularx}{\linewidth}{@{}>{\raggedright\arraybackslash\sffamily\footnotesize\bfseries}p{3.3cm}>{\raggedright\arraybackslash}X>{\raggedright\arraybackslash}X@{}}
\toprule
Criterion & \sffamily\footnotesize\bfseries Version A & \sffamily\footnotesize\bfseries Version B \\
\midrule
1- Opening gives readers a role & \stM{} \sref{1} only names a topic: ``This section discusses\ldots''
               & \stS{} \sref{1} starts from the readers' practice; \sref{4} points to their question: where does my closure fail? \\
2- Names and challenges what readers believe & \stM{} \sref{3} notes the practice (``widely used''), but nothing turns against it
               & \stS{} \sref{1} states the assumption; \sref{2} turns: ``incomplete'' \\
3- Shows what is at stake for readers & \stM{} Nothing is at risk
               & \stP{} \sref{3} states the cost as physics, not as a cost to these readers \\
4- Offers something readers can use & \stM{} The contribution is an activity: ``is analyzed''
               & \stS{} \sref{4}: ``a tool for locating'' \\
5- Sentence subjects are what readers care about & \stP{} Subjects: this section, subgrid-scale models, the Smagorinsky model, the energy transfer. The
                 writer's frame, not the readers' problem
               & \stP{} In \srange{1}{3} the subjects follow the readers (practitioners, the DNS evidence, the model).
                 In \sref{4} the subject is ``this paper''; ``modelers'' would keep the focus on them \\
6- Sentence endings carry the key point & \stP{} \sref{4} ends on method: ``filtered DNS data''
               & \stS{} Endings build: resolved scales $\to$ gaining it $\to$ most likely to fail \\
7- Long sentences pay off early & \stN{} The sentences are short
               & \stS{} \sref{1} and \sref{3} are long but pay off in their first clause \\
8- No sentence needs a reread & \stP{} \sref{3} holds its subject across a 14-word clause before the verb
               & \stS{} No sentence needs a reread \\
9- Nothing slows readers down & \stS{} Easy to read, and just as easy to stop reading
               & \stS{} \emph{A priori test} and \emph{backscatter} are terms these readers use \\
10- Changes what readers think, not only shows knowledge & \stM{} \srange{1}{4} show what the writer knows; none changes what readers think or do
               & \stS{} Each sentence moves readers toward testing their closure \\
\bottomrule
\end{tabularx}
\renewcommand{\arraystretch}{1}
}

{\small \textbf{Neither version is right in the abstract.} Version B works for readers who use eddy-viscosity
closures. For readers who don't, \sref{1} describes someone else's belief and the paragraph loses its hold on
them (G1). B is also longer than A: concision is about the reader's time, not the word count (G4).\par}

% =====================================================================
\clearpage
\section{Worksheet}

\textbf{Reader card}\par
\renewcommand{\arraystretch}{1.6}
\begin{tabularx}{\linewidth}{@{}>{\sffamily\small\bfseries}p{3.4cm}X@{}}
Readers            & \hrulefill \\
Current belief     & \hrulefill \\
What they value    & \hrulefill \\
Function           & \hrulefill \\
Reading conditions & \hrulefill \\
\end{tabularx}

\textbf{Sentence map} (\id{EW-5}, \id{EW-6}, \id{EW-9}, \id{EW-10})\par
\renewcommand{\arraystretch}{1.9}
{\small
\begin{tabularx}{\linewidth}{@{}|c|X|c|X|c|c|X|@{}}
\hline
\sffamily\bfseries S & \sffamily\bfseries Subject (focus) & \sffamily\bfseries Care? & \sffamily\bfseries Last words (stress)
  & \sffamily\bfseries Key? & \sffamily\bfseries C/D & \sffamily\bfseries Friction \\ \hline
1 & & & & & & \\ \hline
2 & & & & & & \\ \hline
3 & & & & & & \\ \hline
4 & & & & & & \\ \hline
5 & & & & & & \\ \hline
6 & & & & & & \\ \hline
\end{tabularx}
}

\textbf{Criteria} \hfill {\small Opening level (\id{EW-1}):\enspace L0 \enspace L1 \enspace L2 \enspace L3}\par
{\small
\begin{tabularx}{\linewidth}{@{}|l|c|X|X|@{}}
\hline
\sffamily\bfseries ID & \sffamily\bfseries Status & \sffamily\bfseries Evidence (S\#) & \sffamily\bfseries Fix move \\ \hline
EW-1  & & & \\ \hline
EW-2  & & & \\ \hline
EW-3  & & & \\ \hline
EW-4  & & & \\ \hline
EW-5  & & & \\ \hline
EW-6  & & & \\ \hline
EW-7  & & & \\ \hline
EW-8  & & & \\ \hline
EW-9  & & & \\ \hline
EW-10 & & & \\ \hline
\end{tabularx}
}
\renewcommand{\arraystretch}{1}

\textbf{Priority fix} (the one criterion that most limits value, and why): \hrulefill\par
\vspace{4pt}\hrulefill

% =====================================================================
\section{Report format for skills and reviewers}

Analyses produced from this handout, by a person or by a skill, use this layout so they can be compared and
cited.

\begin{panel}[breakable=false]
\footnotesize
\begin{verbatim}
Effective-writing report - Self-check | Dissect (someone else's paragraph, G6)
Reader card : readers | current belief | value | function | conditions
              (say which fields were assumed)
Paragraph   : exactly as given, one sentence per line: S1, S2, ... (P1 S1 ...)
Labels      : "1- Opening gives readers a role" to "10- Changes what readers
              think, not only shows knowledge"; plain words throughout, no IDs,
              codes or abbreviations: the report must make sense without this handout
Table       : all ten criteria, in order, in three columns
  SELF-CHECK  | Criterion | Met? | What the paragraph does (citing S#) |
  DISSECT     the same table; no How to fix and no Priority fix
  COVERED     after Handout 1, rows 1- to 4- read "Covered by the lens report"
How to fix  : one line per criterion partly met or not met (self-check)
Priority fix: the one criterion that most limits value, and why (self-check)
Untouched   : numbers, results and technical claims are not changed
\end{verbatim}
\end{panel}

% =====================================================================
\section{Common pitfalls}

\begin{itemize}[leftmargin=1.4em]
  \item Opening with what the text is \emph{about} instead of the readers' situation (\id{EW-1}).
  \item Explaining what you know instead of changing what the readers think (\id{EW-10}).
  \item Writing for an imagined general reader, or ``for sixth graders,'' instead of the actual community and
        what it values \ts{52:00}.
  \item Treating text rules as laws, and their violations as errors (P2, G2).
  \item Assuming readers will read to the end, or that they trust you to make it worth it (P4, P5).
  \item Neutral vocabulary that drains urgency from a real problem (\id{EW-3}).
  \item Asking which version of a paragraph is ``best'' without naming its readers (G1).
\end{itemize}

% =====================================================================
\section{References}
\enlargethispage{2\baselineskip}% keeps the references on the same page

{\small
\renewcommand{\section}[2]{}%  suppress the bibliography's own heading
\begin{thebibliography}{9}\setlength{\itemsep}{0pt}\raggedright

\bibitem{mcenerney2015}
L.~McEnerney, \emph{Leadership Lab: Writing Beyond the Academy}, lecture, University of Chicago,
23 January 2015. Posted on YouTube by UChicago Social Sciences, 18 February 2015:
{\urlstyle{same}\url{https://youtu.be/aFwVf5a3pZM}}.

\bibitem{myerson2000}
R.~B. Myerson, ``Large Poisson games,'' \emph{Journal of Economic Theory}, 94(1), 7--45, 2000.
doi:10.1006/jeth.1998.2453. (Example 1 in the lecture.)

\bibitem{gehlbach2014}
S.~Gehlbach, R.~Myerson and T.~Mylovanov, ``A way forward for Ukraine,'' \emph{The New York Times},
19 March 2014. (Example 2.)

\bibitem{myerson-jcr}
R.~B. Myerson, ``A field manual for the cradle of civilization: Theory of leadership and lessons of Iraq,''
\emph{Journal of Conflict Resolution}, 53(3), 470--482, 2009. (Example 3.)

\bibitem{williams1981}
J.~M. Williams, ``The phenomenology of error,'' \emph{College Composition and Communication}, 32(2),
152--168, 1981. (The article with about 100 planted errors that McEnerney describes; Williams is the
mentor he mentions.)

\bibitem{lincoln1865}
A.~Lincoln, Second Inaugural Address, 4 March 1865. (Final sentence used for the stress example.)

\bibitem{wittgenstein1953}
L.~Wittgenstein, \emph{Philosophical Investigations}, Blackwell, 1953. (Source of ``form of life'' and
``language game.'')

\bibitem{gopen1990}
G.~D. Gopen and J.~A. Swan, ``The science of scientific writing,'' \emph{American Scientist}, 78(6), 550--558,
1990. (Further reading, not named in the lecture: the standard written source for \id{EW-5} and \id{EW-6}.)

\bibitem{williams-style}
J.~M. Williams, \emph{Style: Lessons in Clarity and Grace}, Pearson; many editions, later ones with J.~Bizup.
(Further reading, not named in the lecture.)

\end{thebibliography}

}

\end{document}
